Table~\ref{tab:results-main} compares Red-LGCP with all five baselines; statistical significance is assessed with paired two-sided exact Wilcoxon signed-rank tests over the eight test weeks (four per month; stochastic methods averaged over seeds), with Holm correction across baselines for each metric (App.~\ref{app:significance}). Red-LGCP achieves the lowest MAE in both months and the best Wasserstein distance among the genuinely predictive models (i.e.\ excluding Last Available Poisson, whose near-zero Wasserstein is a trivial artifact of copying real observed counts rather than a sign of good calibration), and both advantages are significant: it is better in at least seven of the eight weeks in every comparison. Most strikingly, its accuracy and F1 far exceed those of Global Cell Mean Poisson and of the covariate-based baselines, which score identically regardless of architecture (GLM, GNN or ConvLSTM). Its F1 is also significantly higher than that of Last Available Poisson, whose accuracy it exceeds in six of eight weeks ($p=0.055$). As the ablation in App.~\ref{app:ablation} confirms, this gap is driven by the zero-inflation classifier, an activity-detection mechanism that none of the baselines has.

ConvLSTM and GNN are the baselines closest to Red-LGCP in a pointwise sense, with the best and second-best RMSE in both months and MAEs close to, though significantly higher than, that of Red-LGCP. This reflects a genuine trade-off rather than a shortcoming, since RMSE weights large individual errors more heavily than MAE, while Red-LGCP's zero-gate is tuned to get the common case (and the activity/inactivity decision) right rather than to avoid the occasional large error on high counts; indeed, none of its RMSE differences from the baselines is statistically significant.
